% -*- coding: utf-8 -*-
% !TEX program = xelatex
\documentclass[plain,noanswer,medmath]{../../template/jnuexam}
\usepackage{../../template/kysx}

\begin{document}


\examtitle{1987}{2}


\loadproblems{1}{1987P1}%载入试卷一


\makepart{填空题}{本题满分15分，每小题3分}%【同试卷一第一题】


\useproblem{1}{1}{1}

\useproblem{1}{1}{2}

\useproblem{1}{1}{3}

\useproblem{1}{1}{4}

\useproblem{1}{1}{5}


\makepart{计算题}{本题满分14分}


\begin{problem}[points=6]
计算定积分$\int_{-2}^2 (|x|+x)\e^{-|x|} \dx$．
\end{problem}


\useproblem[points=8]{1}{2}{0}%【同试卷一第二题】


\makepart{}{本题满分7分}

\begin{problem}
设$z = f(u,x,z)$, $u = xey$，其中$f$具有二阶连续偏导数，求$\frac{\pd^2 z}{\pdx\pd y}$．
\end{problem}


\makepart{}{本题满分8分}

\useproblem{1}{4}{0}%【同试卷一第四题】


\makepart{选择题}{本题满分12分，每小题3分}%【同试卷一第五题】


\useproblem{1}{5}{1}


\useproblem{1}{5}{2}


\useproblem{1}{5}{3}


\useproblem{1}{5}{4}


\makepart{}{本题满分10分}%【同试卷一第六题】

\useproblem{1}{6}{0}


\makepart{}{本题满分10分}%【同试卷一第七题】

\useproblem{1}{7}{0}


\makepart{}{本题满分10分}%【同试卷一第八题】

\useproblem{1}{8}{0}


\makepart{}{本题满分8分}%【同试卷一第九题】

\useproblem{1}{9}{0}


\makepart{}{本题满分6分}

\begin{problem}
设$\lambda_1$,$\lambda_2$为$n$阶方阵$A$的特征值，且$\lambda_1 \ne \lambda_2$，
而$x_1$,$x_2$分别为对应的特征向量，试证明$x_1 + x_2$不是$A$的特征向量．
\end{problem}


\end{document}
